% Generated draft: requires mathematical and source review.
\documentclass[11pt]{article}
\usepackage[margin=25mm]{geometry}
\usepackage{fontspec}
\setmainfont{Latin Modern Roman}
\usepackage{amsmath,amssymb,mathtools,unicode-math}
\setmathfont{Latin Modern Math}
\usepackage{enumitem,needspace,xurl,hyperref}
\setlength{\parindent}{0pt}
\setlength{\parskip}{5pt}
\setlength{\emergencystretch}{3em}
\newcommand{\N}{\mathbb{N}}
\newcommand{\Z}{\mathbb{Z}}
\newcommand{\Q}{\mathbb{Q}}
\newcommand{\R}{\mathbb{R}}
\newcommand{\C}{\mathbb{C}}
\newcommand{\F}{\mathbb{F}}
\newcommand{\A}{\mathbb{A}}
\newcommand{\PP}{\mathbb{P}}
\newcommand{\E}{\mathbb{E}}
\newcommand{\eps}{\varepsilon}
\newcommand{\dd}{\,\mathrm{d}}
\begin{document}
\section*{UnsolvedMath ID: 800307}
\textbf{Source problem code:} catalogue entry
\subsection{Definitions and mathematical statement}
Let \(S^1=\R/\Z\), equipped with the standard circle distance
\[
d(\bar x,\bar y)=\min_{k\in\Z}|x-y-k|.
\]
A homeomorphism \(f:S^1\to S^1\) is Lipschitz with constant at most \(L\) if
\[
d(f(x),f(y))\leq Ld(x,y)\qquad\text{for all }x,y\in S^1.
\]
Let \(G\leq\operatorname{Diff}^{1}(S^1)\) be a finitely generated group, and suppose that \(G\) has no resilient pairs. The supplied record does not define a resilient pair; this hypothesis must therefore retain its source-specific meaning pending verification of the definition.

For a fixed finite generating set \(S\subseteq G\), the requested conclusion is
\[
\forall\varepsilon>0\ \exists h\in\operatorname{Homeo}(S^1)\ \forall s\in S\ \forall x,y\in S^1,\qquad
 d\bigl((h\circ s\circ h^{-1})(x),(h\circ s\circ h^{-1})(y)\bigr)
 \leq e^{\varepsilon}d(x,y).
\]
The question is whether the stated hypotheses imply this conclusion, so that the conjugate group \(hGh^{-1}\) consists of Lipschitz homeomorphisms with the displayed bound on its specified generators. The conjugating homeomorphism may depend on \(\varepsilon\). The record does not specify the generating set or whether it is symmetric; the displayed bound is imposed on the elements of \(S\), with no additional bound on their inverses unless those inverses belong to \(S\). No orientation-preserving restriction is stated in the supplied record. A fully self-contained formulation remains contingent on clarification of the resilient-pair hypothesis and the generating-set convention.
\subsection{Short English statement}
Does every finitely generated group of \(C^1\) circle diffeomorphisms having no resilient pairs admit, for each \(\varepsilon>0\), a topological conjugacy making its generators Lipschitz with constants at most \(e^{\varepsilon}\)? The source-specific meaning of resilient pairs and the generating-set convention require clarification.
\subsection{Sources}
\begin{itemize}
\item Supplied local UnsolvedMath registry record, permanent ID \(800307\), source code \texttt{catalogue entry}, titled \emph{Group Actions on 1-Manifolds --- Question 24}. The registry attributes the question to Navas, with proposed year 2017, and labels it open. This is a supplied status label, not an independently verified literature assessment.
\item Navas, \emph{Group Actions on 1-Manifolds: A List of Very Concrete Open Questions} (2017), \url{https://arxiv.org/abs/1712.06462}. The supplied provenance identifies the item as ``Question 24 (\#23)'', records access on 2026-10-16, and gives extraction method \texttt{arxiv\_latex}. These provenance details are supplied metadata; the source has not been independently consulted for this draft.
\end{itemize}
\end{document}
